% Converted from: paper/sections/02_platform_design_and_implementation.md
\section{Platform Design and Implementation}

\subsection{Mission Scenario and Engineering Requirements}

The platform studied in this work is developed for underwater corner-reflector deployment, where the reflector acts as a passive marker and the carrier must rapidly achieve an operational posture after release. In the considered scenario, the platform is packaged in a compact, near-horizontal configuration and released with nonzero initial velocity. The mission-critical requirement is to complete the horizontal-launch to vertical-stabilization transition reliably, because the payload can be regarded as operationally effective only after the platform reaches and maintains a near-vertical working posture.

This requirement imposes coupled constraints on the platform: the mechanical architecture and buoyancy layout must support a stable working posture, the transient response must avoid excessive overshoot and residual rate, and the design must remain compatible with practical deployment and packaging constraints. These constraints motivate treating the platform design and fabrication as a primary contribution and motivate an accompanying transition-focused model and evidence chain.

\subsection{Realized Platform Architecture and Implementation}

The realized platform is a mission-oriented self-suspending system whose architecture is defined by a rigid frame, a prescribed buoyancy-center offset relative to the mass distribution, and structural features that induce inner-water coupling during large-attitude transition. The platform is designed such that the vertical working posture is a stable equilibrium under hydrostatic restoring, while the transient motion remains compatible with deployment constraints and the available actuation strategy.

A high-level representation of the hardware architecture and the deployment workflow is summarized in Fig.~\ref{fig:mission-architecture}. Detailed mission-sensitive operational information is intentionally abstracted, but the physical and experimental assumptions that affect the transition dynamics are stated explicitly so that the modeling and validation results remain interpretable and reproducible.

\subsection{Mass--Buoyancy Arrangement and Self-Suspending Mechanism}

The self-suspending behavior is achieved through the relative placement of the centers of buoyancy and gravity, which generates a restoring moment that drives the platform toward the working posture. For the nominal configuration used throughout this paper, the baseline physical parameters are summarized in Table~\ref{tab:baseline-params}.

\begin{table}[htbp]
  \centering
  \caption{Baseline physical parameters of the realized platform.}
  \label{tab:baseline-params}
  \begin{tabular}{llrl}
    \toprule
    Parameter & Symbol & Value & Unit \\
    \midrule
    Water density & $\rho$ & 997.561 & kg/m$^3$ \\
    Reference area & $A_{ref}$ & 0.0056745 & m$^2$ \\
    Reference length & $L_{ref}$ & 0.625 & m \\
    Dry mass & $m_{dry}$ & 2.55 & kg \\
    Wet mass & $m_{wet}$ & 2.76 & kg \\
    Inner-water equivalent mass & $m_{water,inner}$ & 0.21 & kg \\
    Pitch inertia & $I_{yy}$ & 0.05741 & kg$\cdot$m$^2$ \\
    Inner-water equivalent pitch inertia & $I_{water,inner}$ & 0.01119 & kg$\cdot$m$^2$ \\
    Buoyancy equivalent mass & $B_{mass}$ & 2.55 & kg \\
    Buoyancy offset (body frame) & $x_b$ & 0.02535 & m \\
    Buoyancy offset (body frame) & $z_b$ & 0.0 & m \\
    \bottomrule
  \end{tabular}
\end{table}

For model closure and later sensitivity analysis, buoyancy is represented through an equivalent buoyancy mass $B_{mass}$ and a nominal body-frame offset $(x_b, z_b)$. The restoring moment generated by the buoyancy offset provides the dominant passive mechanism for posture recovery during the transition, while damping and inner-water coupling shape the transient response.

\subsection{Deployment Workflow and Analyzed Phase Definition}

This paper targets a single mission-critical phase: horizontal release followed by passive transition toward vertical stabilization. The stabilized working posture lies near $\theta=90^\circ$ under the sign convention used in this study. The analysis focuses on the transient evolution from the release state to the stabilized posture and excludes downstream tasks such as depth regulation or long-horizon trajectory management.

The phase definition is aligned with the deployment workflow in Fig.~\ref{fig:mission-architecture} and provides the boundary conditions for the experimental program (Section~3), the reduced-order model (Section~4), and the validation and design analyses (Sections~6--7).

\subsection{Coordinate Conventions and State Interface}

A fixed body-axis convention is used throughout the repository and manuscript to avoid sign inconsistencies across CFD post-processing, parameter identification, and time-domain simulation. The body axes are defined as $x_b$ forward and $z_b$ downward. Positive pitch is nose-up, and pitch rate is $q=\dot{\theta}$ about $+y_b$. The state interfaces are $\nu=[u,w,q]^T$ and $\eta=[\theta]$, and the angle of attack for hydrodynamic lookup is defined by $\alpha=\mathrm{atan2}(w,u)$. These conventions are frozen in the implementation and should be treated as part of the model contract.

\begin{figure}[htbp]
  \centering
  % TODO: Insert Fig. 2 here (see paper/figures/FIGURE_REQUIREMENTS.md).
  \caption{Body-axis definition, sign convention, and state interface ($u$, $w$, $q$, $\theta$).}
  \label{fig:conventions}
\end{figure}

\subsection{System Boundary and Assumptions}

Three assumptions delimit the claims made in the remainder of the paper. First, mission context is retained at the level required to motivate the platform and the transition-phase objective, but mission-sensitive details are abstracted. Second, all conclusions are intentionally restricted to the current platform geometry, mass--buoyancy configuration, and actuation strategy. Third, CFD is used as hydrodynamic closure and mechanism-support evidence, not as a standalone novelty claim independent of the platform--model--validation workflow.
